\section{The unbounded binary runtime proof}
\label{sec:runtime}

The complete final-table check proves 8,386 guarded interfaces using
\textbf{2,255,513 symbolic execution rules}. All 476 literal transition
rows occur in these derivations. The minimum constant part of a
certified instruction duration is \textbf{13}. These figures count
proof construction. They do not estimate the machine's ultimate search
time, and the machine's nonhalting is not inferred from a long test run.

\subsection{The quantified configuration family}

Define a unary register block and its transient form by
\[
 B(v)=0\,1^{v+1},\qquad C(v)=1\,0\,1^v
 \quad(v\in\Nat).
\]
Both have length $v+2$. For a finite list $A=(a_1,\ldots,a_k)$,
write $B(A)=B(a_1)\cdots B(a_k)$, and similarly for $C(A)$.
The primary registers are permuted physically by
\[
 \pi=(6,2,0,4,3,7,8,5,1).
\]
That is, logical register $i$ occupies physical block $\pi(i)$.
The physical block vector for the initial logical seed is
$(0,0,0,3,0,0,3,0,0)$.

A clean instruction frame at inherited dispatcher 96 has a one-bit
marker immediately left of the head, a path word followed by zeros,
a one-bit header at distance $64+g$ from the head, and then the
blocks of the nine primary registers followed by a nonempty spare
list. The exterior is zero. The integer $g$ is any nonnegative
additional gap. The source has maximum path length 48, so the path
fits before the fixed minimum header distance 64.

Each symbolic proof independently quantifies:
\begin{itemize}
\item all nine primary values $r_0,\ldots,r_8\in\Nat$;
\item the additional gap $g\in\Nat$;
\item the first spare value $a\in\Nat$;
\item an arbitrary finite remaining spare list $A$.
\end{itemize}
The complete spare suffix is $B(a)B(A)$, and is therefore nonempty
even when $A$ is empty. A decrement case is split into its zero
case and its positive case $r_i=u+1$, with $u\in\Nat$.
No finite bound is placed on these parameters or on the length of
the remaining list.

\subsection{Exact source effects}

The source operations are increment (\term{CI}), decrement-if-positive
(\term{CD}), break (\term{CB}), continue (\term{CC}), and halt
(\term{CH}). Paths are read most-significant bit first with trailing
zero padding. After a leaf is selected, the path is extended or
truncated to the leaf's depth before path arithmetic. This matches
the included \file{RegisterPath.lean} semantics.

Let $\operatorname{inc}(p)$ be fixed-width binary increment, with
all-ones rollover, and let $\operatorname{strip}_b^k(p)$ remove the
suffix through the $k$th occurrence of bit $b$ counted from the right,
or all of $p$ if there are too few. Then:
\begin{center}
\begin{tabular}{@{}L{0.16\textwidth}L{0.32\textwidth}L{0.42\textwidth}@{}}
\toprule
Operation & Logical effect & Successor pointer\\
\midrule
$\term{CI}(i)$ & $r_i\gets r_i+1$ & $\operatorname{inc}(p)$\\
$\term{CD}(i)$, positive & $r_i\gets r_i-1$ & $\operatorname{inc}(p)$\\
$\term{CD}(i)$, zero & Registers unchanged &
$\operatorname{inc}(\operatorname{strip}_0^1(p))$\\
$\term{CB}(k)$ & Registers unchanged &
$\operatorname{inc}(\operatorname{strip}_0^k(p))$\\
$\term{CC}(k)$ & Registers unchanged &
$\operatorname{strip}_1^k(p)$\\
$\term{CH}$ & Literal binary halt & No successor\\
\bottomrule
\end{tabular}
\end{center}

Increment appends one spare block of value 1. Other operations
preserve all existing spare values. A successful decrement increases
the additional gap by one; binary path rollover can also increase
the gap and shift the head left by one. The target frame includes
that exact displacement. In particular, the next gap is nonnegative
and the spare suffix remains nonempty.

\subsection{Three sound execution rules}

The proof engine represents each homogeneous run length by an affine
form
\[
 c+\alpha g+\sum_{i=0}^{8}\beta_i r_i+\gamma a+\delta W,
\]
where $W$ denotes the width of the remaining spare list when needed.
There are thirteen coefficient slots including the constant. It only
uses a run as uniformly nonempty when the constant is at least one;
all variable coefficients must be nonnegative. It refuses an
unjustified read into an unknown or possibly empty context.

\begin{lemma}[Literal transition rule]
If the current cell is represented by a run with a positive constant
length and known bit, the engine's one-transition rule has exactly
the effect of the corresponding literal table row, for every
assignment of natural parameters.
\end{lemma}
\begin{proof}
Writing changes the current bit and moving changes the two sides of
the zipper by one cell. The affine run is shortened by exactly one.
Its remaining coefficients stay nonnegative. The next state, direction,
duration increment, and displacement are read from that literal row.
No information about any other cell is altered. This is the defining
Turing-machine transition equation.
\end{proof}

\begin{lemma}[Homogeneous-run induction]
Suppose a row, reading $b$, writes $w$, moves in the current scan
direction, and returns to the same state. Traversing a run $b^n$
has the effect of replacing it by $w^n$, moving exactly $n$ cells
in that direction, and retaining the state, for every $n\in\Nat$.
\end{lemma}
\begin{proof}
For $n=0$ the statement is the identity. For $n+1$, the first literal
transition writes $w$ and exposes a run $b^n$ in the same state.
Apply the induction hypothesis. Durations and displacements add.
This also explains why accelerating a symbolic run whose value can
be zero is sound: the zero instance takes zero literal transitions.
\end{proof}

The only arbitrary-list accelerations are founded on four one-block
schemas. Their names identify entry controls, while their validity
is rederived from the \emph{actual supplied table}:
\begin{center}
\begin{tabular}{@{}rccc@{}}
\toprule
Entry state & Direction & Block effect & Exact duration\\
\midrule
73 & Right & $B(v)\mapsto B(v)$ & $v+2$\\
95 & Right & $B(v)\mapsto C(v)$ & $v+2$\\
48 & Left & $B(v)\mapsto B(v)$ & $v+2$\\
112 & Left & $C(v)\mapsto B(v)$ & $v+2$\\
\bottomrule
\end{tabular}
\end{center}
For a leftward scan, the block is encountered in reverse order.
The proofs use an arbitrary opaque context on each side. Reaching an
opaque context prematurely is a failure, so a one-block proof
cannot silently assume a bit outside its stated footprint.

\begin{lemma}[Finite-list induction]
\label{lem:list}
Each verified one-block schema above extends to every finite list
of blocks, with the same entry/exit state and total duration equal
to the sum of their lengths.
\end{lemma}
\begin{proof}
The empty list requires no transition. For a nonempty list, apply
the one-block result to its first encountered block, taking the
remainder as the arbitrary exterior context. Apply the induction
hypothesis to the remaining list. The state at the join is exactly
the next schema's entry state. The tape transformation and displacement
therefore concatenate, including the leftward reversal convention.
\end{proof}

Normalization also uses the word identity
\[
 C(A)\,1=1\,B(A),
\]
and its reversal. For one block,
$C(v)1=101^{v+1}=1B(v)$. The empty-list case and induction on the
list prove the general identity. Adjacent equal-symbol runs are
joined, and an infinite zero exterior absorbs a finite adjoining
zero run. Each normalization is an exact word equality.

\subsection{What acceptance of a case proves}

For a nonhalting operation, acceptance requires equality of the
entire symbolic output tape with the prescribed successor frame,
equality of state and orientation, and the exact affine head
displacement. It also requires a duration with nonnegative
coefficients and strictly positive constant part.
For the halt instruction, the literal halt state must actually
be reached after a positive finite segment. An intermediate
unmatched halt would have no next transition and cause failure.

The case driver checks every leaf and both outcomes of every
decrement, not only combinations observed during execution.
There are 1,931 decrement leaves, giving
$6,455+1,931=8,386$ cases. The finite cap of 4,000 symbolic rules
per derivation limits the effort spent constructing a certificate;
it does not limit any parameter of the theorem or the number of
literal transitions represented by a derived rule.

For example, an instruction segment whose duration is
$c+\beta r_i+\delta W$ may require arbitrarily many literal
transitions. The same finite derivation proves its correctness for
every $r_i$ and every finite spare list. This is the reason that a
finite successful case census can establish an unbounded runtime
theorem.

\begin{proposition}[Final-table instruction interface]
\label{prop:runtime}
Every guarded leaf of the supplied boundary source is implemented
by the exact 238-state table on the complete quantified frame
family above. Each nonhalting leaf reaches its prescribed successor
frame in finite positive time; the halt leaf reaches the designated
halt outcome.
\end{proposition}
\begin{proof}
The attributed engine proves the four one-block schemas and then
constructs each case using only the three rules just justified and
exact word normalization. The complete driver accepts all 8,386
cases and checks the literal row census, source completeness, exact
source hash, positive clocks, and expected aggregate counts.
Induction over each recorded derivation proves its universally
quantified interface. Lemma~\ref{lem:list} discharges the only
arbitrary-list rule. The results and every affine case record are
included in \file{evidence/}.
\end{proof}

\subsection{Exact initialization}
\label{sec:boot}

The new table, from the prescribed blank start, reaches this frame
after exactly 409,835 transitions:
\begin{center}
\begin{tabular}{@{}L{0.38\textwidth}L{0.51\textwidth}@{}}
\toprule
Frame component & Value\\
\midrule
State; head & Canonical 84; $-5$\\
Marker; header & $-6$; $176$\\
Header distance; extra gap & $181$; $117$\\
Logical primary values & $(3,0,0,0,3,0,0,0,0)$\\
Physical primary values & $(0,0,0,3,0,0,3,0,0)$\\
Spare suffix & 167 blocks, each of value 1\\
One-cell count & 351\\
Other tape positions & All zero\\
\bottomrule
\end{tabular}
\end{center}

The checker runs a sparse finite-support interpreter, an independent
dense-array interpreter, and a dense run of the canonical renumbering.
It compares their complete states, heads, and one-cell supports,
then parses the actual unary blocks to recover the displayed
registers and spares. The new complete tape also equals the pinned
calibrated startup frame. This is a finite proof of a specific
initialization, not an argument for infinite nonhalting.

The source/frame lemma in Section~\ref{sec:source} now closes the
runtime invariant from this initial configuration: every successor
frame is literally the input frame for the next source leaf.
Consequently Proposition~\ref{prop:runtime} applies at every future
source step, with no hidden bound on the growing computation.
